% GENERATED FILE. Edit canonical lessons, then run npm run generate:books.
% canonical-source-hash: fnv1a64:f7cb819ae88cd69b
% canonical-lessons: LA-C302-speculator, LA-C302-explorator, LA-C302-exploro, LA-C302-prodo, LA-C302-patefacio

\chapter{Speculātōrēs}
\label{ch:la-speculatores}
\hlchaptermodality{302}

\begin{chapteropening}
\textbf{By the end of this chapter:} \emph{I can tell how secrets came out: a spy, a scout, to search out, to betray and to disclose.}

Complete the last lesson of chapter 302: I can tell how secrets came out: a spy, a scout, to search out, to betray and to disclose.
\end{chapteropening}

\section[speculātor, speculātōris]{speculātor, speculātōris — a spy, a lookout}

\label{lesson:LA-C302-speculator}



\begin{quote}
Before the new one: say the Latin for to grow stronger, then the Latin for to assert.
\end{quote}

\subsection*{You'll want to know: speculātor, speculātōris}
\textbf{speculātor, speculātōris} — \textquotedblleft{}a spy, a lookout\textquotedblright{}.

\begin{quote}
Speculātor noctū in castra hostium vēnit.
\end{quote}

\begin{quote}
\emph{A spy came into the enemy camp by night.}
\end{quote}

Masculine, 3rd declension. A soldier sent ahead to watch. English \emph{speculator} now means someone who gambles on prices.

\begin{cousinweb}
\textbf{speculātor} is built on \emph{speculor}, \textquotedblleft{}to keep watch\textquotedblright{}, from \emph{specula}, \textquotedblleft{}a lookout\textquotedblright{}, which belongs with \emph{speciō}, \textquotedblleft{}to look\textquotedblright{} (Lewis \& Short). English \emph{speculate} is from the same verb.
\end{cousinweb}

\begin{grammarlens}[title={What you've built}]
One more word for telling, explaining and describing at B1.
\end{grammarlens}

\subsection*{Your turn}
\begin{itemize}
\raggedright
  \item \emph{Say it:} \emph{speculātor, speculātōris}
  \item \emph{Say it:} \emph{speculātor, speculātōris}, once more
  \item \emph{Say it:} a sentence of your own with \emph{speculātor, speculātōris}
  \item \emph{Recall:} say the Latin for to grow stronger, then the Latin for to assert, then say \emph{speculātor, speculātōris} again
\end{itemize}

\begin{tcolorbox}[breakable,colback=teal!4,colframe=teal!35!black,title={Before you move on}]
What does speculātor, speculātōris mean? (\textquotedblleft{}A spy, a lookout\textquotedblright{}.) Say it once more.
\end{tcolorbox}
\section[explōrātor, explōrātōris]{explōrātor, explōrātōris — a scout}

\label{lesson:LA-C302-explorator}



\begin{quote}
Before the new one: say the Latin for to assert, then the Latin for a spy.
\end{quote}

\subsection*{You'll want to know: explōrātor, explōrātōris}
\textbf{explōrātor, explōrātōris} — \textquotedblleft{}a scout\textquotedblright{}.

\begin{quote}
Explōrātōrēs nūntiāvērunt hostēs prope flūmen esse.
\end{quote}

\begin{quote}
\emph{The scouts reported that the enemy were near the river.}
\end{quote}

Masculine, 3rd declension. In Caesar the \emph{explōrātōrēs} went ahead of the army.

\begin{cousinweb}
\textbf{explōrātor} is built on \emph{explōrō}, \textquotedblleft{}to search out\textquotedblright{} (Lewis \& Short), with the ending \emph{-tor} for the one who does it.
\end{cousinweb}

\begin{grammarlens}[title={What you've built}]
One more word for telling, explaining and describing at B1.
\end{grammarlens}

\subsection*{Your turn}
\begin{itemize}
\raggedright
  \item \emph{Say it:} \emph{explōrātor, explōrātōris}
  \item \emph{Say it:} \emph{explōrātor, explōrātōris}, once more
  \item \emph{Say it:} a sentence of your own with \emph{explōrātor, explōrātōris}
  \item \emph{Recall:} say the Latin for to assert, then the Latin for a spy, then say \emph{explōrātor, explōrātōris} again
\end{itemize}

\begin{tcolorbox}[breakable,colback=teal!4,colframe=teal!35!black,title={Before you move on}]
What does explōrātor, explōrātōris mean? (\textquotedblleft{}A scout\textquotedblright{}.) Say it once more.
\end{tcolorbox}
\section[explōrō, explōrāre]{explōrō, explōrāre — to search out, to explore}

\label{lesson:LA-C302-exploro}



\begin{quote}
Before the new one: say the Latin for a spy, then the Latin for a scout.
\end{quote}

\subsection*{You'll want to know: explōrō, explōrāre}
\textbf{explōrō, explōrāre} — \textquotedblleft{}to search out, to explore\textquotedblright{}.

\begin{quote}
Mīlitēs silvam explōrāvērunt, sed nēminem invēnērunt.
\end{quote}

\begin{quote}
\emph{The soldiers searched the forest, but found nobody.}
\end{quote}

\emph{explōrō, explōrāre, explōrāvī, explōrātum}. English \emph{explore}.

\begin{cousinweb}
\textbf{explōrō} is usually traced to \emph{ex} and \emph{plōrō}, \textquotedblleft{}to cry out\textquotedblright{}, perhaps first the shout of hunters driving game out of cover; the details are not settled. English \emph{explore} comes from it.
\end{cousinweb}

\begin{grammarlens}[title={What you've built}]
One more word for telling, explaining and describing at B1.
\end{grammarlens}

\subsection*{Your turn}
\begin{itemize}
\raggedright
  \item \emph{Say it:} \emph{explōrō, explōrāre}
  \item \emph{Say it:} \emph{explōrō, explōrāre}, once more
  \item \emph{Say it:} a sentence of your own with \emph{explōrō, explōrāre}
  \item \emph{Recall:} say the Latin for a spy, then the Latin for a scout, then say \emph{explōrō, explōrāre} again
\end{itemize}

\begin{tcolorbox}[breakable,colback=teal!4,colframe=teal!35!black,title={Before you move on}]
What does explōrō, explōrāre mean? (\textquotedblleft{}To search out, to explore\textquotedblright{}.) Say it once more.
\end{tcolorbox}
\section[prōdō, prōdere]{prōdō, prōdere — to betray; to hand down}

\label{lesson:LA-C302-prodo}



\begin{quote}
Before the new one: say the Latin for a scout, then the Latin for to search out.
\end{quote}

\subsection*{You'll want to know: prōdō, prōdere}
\textbf{prōdō, prōdere} — \textquotedblleft{}to betray; to hand down\textquotedblright{}.

\begin{quote}
Servus dominum hostibus prōdidit.
\end{quote}

\begin{quote}
\emph{The slave betrayed his master to the enemy.}
\end{quote}

\emph{prōdō, prōdere, prōdidī, prōditum}: to give forth. Also 'to hand down': \emph{memoriae prōdere}, to record for later ages.

\begin{cousinweb}
\textbf{prōdō} is built on \emph{dō}, \textquotedblleft{}to give\textquotedblright{}, with \emph{prō}, \textquotedblleft{}forth\textquotedblright{} (Lewis \& Short): to give forth, to give up, and so to betray.
\end{cousinweb}

\begin{grammarlens}[title={What you've built}]
One more word for telling, explaining and describing at B1.
\end{grammarlens}

\subsection*{Your turn}
\begin{itemize}
\raggedright
  \item \emph{Say it:} \emph{prōdō, prōdere}
  \item \emph{Say it:} \emph{prōdō, prōdere}, once more
  \item \emph{Say it:} a sentence of your own with \emph{prōdō, prōdere}
  \item \emph{Recall:} say the Latin for a scout, then the Latin for to search out, then say \emph{prōdō, prōdere} again
\end{itemize}

\begin{tcolorbox}[breakable,colback=teal!4,colframe=teal!35!black,title={Before you move on}]
What does prōdō, prōdere mean? (\textquotedblleft{}To betray; to hand down\textquotedblright{}.) Say it once more.
\end{tcolorbox}
\section[patefaciō, patefacere]{patefaciō, patefacere — to open up; to disclose}

\label{lesson:LA-C302-patefacio}



\begin{quote}
Before the new one: say the Latin for to search out, then the Latin for to betray.
\end{quote}

\subsection*{You'll want to know: patefaciō, patefacere}
\textbf{patefaciō, patefacere} — \textquotedblleft{}to open up; to disclose\textquotedblright{}.

\begin{quote}
Ūna epistula tōtum cōnsilium patefēcit.
\end{quote}

\begin{quote}
\emph{A single letter disclosed the whole plan.}
\end{quote}

\emph{patefaciō, patefacere, patefēcī, patefactum}. Cicero uses it of plots brought to light.

\begin{cousinweb}
\textbf{patefaciō} joins \emph{pateō}, \textquotedblleft{}to be open\textquotedblright{}, and \emph{faciō}, \textquotedblleft{}to make\textquotedblright{} (EtymDB): to make open.
\end{cousinweb}

\begin{grammarlens}[title={What you've built}]
One more word for telling, explaining and describing at B1.
\end{grammarlens}

\subsection*{Your turn}
\begin{itemize}
\raggedright
  \item \emph{Say it:} \emph{patefaciō, patefacere}
  \item \emph{Say it:} \emph{patefaciō, patefacere}, once more
  \item \emph{Say it:} a sentence of your own with \emph{patefaciō, patefacere}
  \item \emph{Recall:} say the Latin for to search out, then the Latin for to betray, then say \emph{patefaciō, patefacere} again
\end{itemize}

\begin{tcolorbox}[breakable,colback=teal!4,colframe=teal!35!black,title={Before you move on}]
What does patefaciō, patefacere mean? (\textquotedblleft{}To open up; to disclose\textquotedblright{}.) Say it once more.
\end{tcolorbox}
